% Generated draft: requires mathematical and source review.
\documentclass[11pt]{article}
\usepackage[margin=25mm]{geometry}
\usepackage{fontspec}
\setmainfont{Latin Modern Roman}
\usepackage{amsmath,amssymb,mathtools,unicode-math}
\setmathfont{Latin Modern Math}
\usepackage{enumitem,needspace,xurl,hyperref}
\setlength{\parindent}{0pt}
\setlength{\parskip}{5pt}
\setlength{\emergencystretch}{3em}
\newcommand{\N}{\mathbb{N}}
\newcommand{\Z}{\mathbb{Z}}
\newcommand{\Q}{\mathbb{Q}}
\newcommand{\R}{\mathbb{R}}
\newcommand{\C}{\mathbb{C}}
\newcommand{\F}{\mathbb{F}}
\newcommand{\A}{\mathbb{A}}
\newcommand{\PP}{\mathbb{P}}
\newcommand{\E}{\mathbb{E}}
\newcommand{\eps}{\varepsilon}
\newcommand{\dd}{\,\mathrm{d}}
\begin{document}
\section*{UnsolvedMath ID: 30004602}
\textbf{Source problem code:} catalogue entry
\subsection{Definitions and mathematical statement}
The supplied record does not specify the field hypotheses or define the Jacobian ideal and the third weak Lefschetz property. The following gives an explicit conventional interpretation; these conventions require confirmation against the cited sources.

Let \(K\) be the intended field and let \(S=K[x_1,x_2,x_3]\), graded by \(\deg x_j=1\). A central hyperplane arrangement \(\mathcal A\) in \(K^3\) is a finite set of distinct linear hyperplanes through the origin. Choose a nonzero linear equation \(\alpha_H\in S_1\) for each \(H\in\mathcal A\), and put
\[
Q_{\mathcal A}=\prod_{H\in\mathcal A}\alpha_H.
\]
Under the gradient-ideal convention, define
\[
J_{\mathrm{grad}}(\mathcal A)=\left(\frac{\partial Q_{\mathcal A}}{\partial x_1},\frac{\partial Q_{\mathcal A}}{\partial x_2},\frac{\partial Q_{\mathcal A}}{\partial x_3}\right)\subseteq S.
\]
Another possible convention includes \(Q_{\mathcal A}\) among the generators. Euler's identity identifies these two ideals when \(\deg Q_{\mathcal A}\) is invertible in \(K\); the record supplies no hypothesis ensuring this.

For a graded quotient \(B\) of \(S\), interpret \(3\)-WLP as follows: there exist \(\ell_1,\ell_2,\ell_3\in S_1\) such that, writing
\[
B^{(0)}=B,\qquad B^{(i)}=B^{(i-1)}/\ell_iB^{(i-1)}\quad(1\leq i\leq3),
\]
the multiplication maps satisfy
\[
\operatorname{rank}_K\!\left(\times\ell_i:B^{(i-1)}_d\longrightarrow B^{(i-1)}_{d+1}\right)
=\min\!\left\{\dim_K B^{(i-1)}_d,\dim_K B^{(i-1)}_{d+1}\right\}
\quad(1\leq i\leq3,\ d\geq0).
\]
Here each \(\ell_i\) acts through its image in the indicated quotient. This definition imposes no Artinian hypothesis; whether this is the intended convention is unresolved.

\textbf{Conjectural target under these conventions.} For every central hyperplane arrangement \(\mathcal A\) in \(K^3\), the quotient \(S/J_{\mathrm{grad}}(\mathcal A)\) has \(3\)-WLP. The source writes \(S/J(\mathcal A)\); its precise ideal convention, field restrictions, and Lefschetz convention must be established before this interpretation can be treated as an unqualified formulation of the source conjecture.
\subsection{Short English statement}
The conjecture asks whether every central hyperplane arrangement in three variables has a Jacobian quotient admitting three successive linear multiplication operators of maximal rank in every degree. The field hypotheses and the precise Jacobian and Lefschetz conventions are absent from the supplied record, so the explicit formulation above is provisional.
\subsection{Sources}
\begin{itemize}
\item Local registry: \texttt{data/unsolvedmath/problems.json}, permanent UnsolvedMath ID \(30004602\), source problem code \texttt{catalogue entry}. The record gives proposed year 2021 and names no individual proposer.
\item Supplied report citation: \emph{Logarithmic Vector Fields and Freeness of Divisors and Arrangements: New perspectives and applications} (2022), Oberwolfach Reports, pp.\ 225--300, DOI 10.4171/owr/2021/5. \url{https://doi.org/10.4171/owr/2021/5}. The supplied literature list labels this reference ``Oberwolfach Report 5/2021.''
\item Supplied literature reference: Palezzato--Torielli, \emph{Journal of Algebra}. \url{https://eprints.lib.hokudai.ac.jp/dspace/bitstream/2115/91057/1/J.%20Algebra%20590_215-233.pdf}. No article title or fuller bibliographic citation is supplied.
\item Literature assessment dated exactly \texttt{2026-08-22}: the supplied curation labels the conjecture open and unresolved, reports that the workshop report mentions only an intended proof, and reports failures in higher-dimensional examples. This is supplied dated curation, not an independently verified assessment; time-sensitive status and cited claims require independent verification.
\item The registry labels the statement \texttt{corrected\_verified} and says it was checked against the official report. This is a supplied verification claim, not verification performed for this draft. The background describes the collection as drawn from Oberwolfach workshop reports and their open-problem discussions.
\end{itemize}
\end{document}
